\chapter{IMPLEMENTACIÓN}
\label{ch:implementacion}

Este capítulo consolida la implementación del sistema de navegación autónoma, desde el hardware hasta la arquitectura de software. Se presentan los componentes del prototipo diferencial y sus parámetros físicos, el gemelo digital en Gazebo, la construcción físico-electrónica y su integración con \gls{ROS2} Foxy, el firmware de bajo nivel en \gls{Arduino} Nano y, finalmente, la arquitectura de control y navegación que integra \texttt{ros2\_control}, \texttt{slam\_toolbox}, Nav2 y \texttt{explore\_lite}. Los detalles operativos, configuraciones extendidas y evidencia de verificación se trasladan a los anexos correspondientes.

%=====================================================================
\section{Componentes y parámetros del prototipo}
\label{ch:componentes}

El prototipo corresponde a una plataforma diferencial de bajo costo para interiores: controla el desplazamiento mediante la diferencia de velocidad entre las ruedas motrices izquierda y derecha, lo que facilita la implementación del modelo cinemático desarrollado en el Capítulo~\ref{marco_teorico}. La Tabla~\ref{tab:componentes_principales} resume los componentes empleados y su función dentro del sistema.

\begin{table}[H]
\centering
\caption{Componentes principales del prototipo.}
\label{tab:componentes_principales}
\small
\begin{tabular}{p{3.2cm} p{4.5cm} p{5.2cm}}
\toprule
\textbf{Subsistema} & \textbf{Componente} & \textbf{Función principal} \\
\midrule
Locomoción & Motores DC JGA25-370 con ruedas de goma & Generar el movimiento diferencial del robot. \\
Odometría & Encoders de cuadratura Hall & Medir el desplazamiento de cada rueda. \\
Potencia de motores & Driver L298N & Accionar los motores mediante señales \gls{PWM}. \\
Control de bajo nivel & \gls{Arduino} Nano & Ejecutar lectura de encoders y control PID de velocidad. \\
Cómputo embebido & Raspberry Pi 4 & Ejecutar la interfaz de hardware, adquirir el \gls{LIDAR} y comunicarse con el \gls{Arduino} y la estación de trabajo. \\
Cómputo de alto nivel & Estación de trabajo & Ejecutar \texttt{slam\_toolbox}, Nav2, exploración autónoma y RViz2. \\
Percepción & RPLidar A1 & Obtener barridos láser 2D para mapeo y evasión de obstáculos. \\
Energía & Batería Li-ion 12 V y regulador buck 5 V & Alimentar motores, lógica y sensores. \\
Estructura & Chasis impreso en 3D & Integrar mecánicamente los componentes del robot. \\
\bottomrule
\end{tabular}
\end{table}

\subsection{Actuación, odometría y control de bajo nivel}

La actuación se realiza mediante dos motorreductores JGA25-370 de 12 V, seleccionados por su disponibilidad, bajo costo y torque suficiente para una plataforma diferencial de interiores. Cada motor incorpora un encoder incremental de cuadratura, utilizado para estimar velocidad y desplazamiento de rueda.

La resolución teórica del encoder en el eje de salida se obtiene a partir de la resolución del sensor Hall en el eje del motor, la decodificación en cuadratura y la relación de reducción:

\begin{equation}
N_e^{\mathrm{teo}} = N_{\mathrm{Hall}} \cdot c_q \cdot N_g
\label{eq:resolucion_encoder}
\end{equation}

donde $N_e^{\mathrm{teo}}$ es la resolución teórica del encoder [ticks/vuelta del eje de salida], $N_{\mathrm{Hall}}=11$ el número de pulsos por vuelta del eje del motor, $c_q$ el factor de decodificación en cuadratura y $N_g$ la relación de reducción; sus valores se consolidan en la Tabla~\ref{tab:parametros_modelo}. Las mediciones experimentales, detalladas en la Sec.~\ref{sec:cpr_experimental}, sustentan el valor efectivo $N_e$ de la Tabla~\ref{tab:parametros_modelo}, utilizado de forma unificada por el firmware y la odometría.

El control de bajo nivel se ejecuta en un \gls{Arduino} Nano, encargado de leer los encoders, calcular el PID de velocidad de cada rueda y enviar señales \gls{PWM} al driver L298N. Este último introduce una zona muerta y caída de tensión que deben compensarse mediante la sintonización experimental del controlador.

\subsection{Cómputo, percepción e integración con ROS 2}

La Raspberry Pi 4 ejecuta Ubuntu 20.04 y \gls{ROS2} Foxy, actuando como computador principal del robot. Su función es coordinar la comunicación con el \gls{Arduino}, adquirir y publicar los datos del sensor láser, y ejecutar la interfaz de hardware \texttt{DiffDriveArduino} y el controlador diferencial \texttt{diff\_drive\_controller}; el PID de bajo nivel permanece en el \gls{Arduino}, mientras que SLAM, Nav2 y exploración se ejecutan en la estación de trabajo.

El sensor principal de percepción es el RPLidar A1, un escáner láser 2D de 360 grados que publica barridos utilizados por SLAM y por los mapas de costo de Nav2. Su selección y limitaciones se justifican en la Sec.~\ref{sec:sensores}.

\subsection{Alimentación y fabricación}

El sistema se alimenta desde una batería Li-ion de 12 V; la organización de los dominios de potencia y lógica se detalla en la Sec.~\ref{chap:robot_real}. El chasis y las cubiertas fueron fabricados mediante impresión 3D FDM en PLA+, tomando como referencia una plataforma diferencial comercial \citep{BECChassisLong}. El diseño considera espacio para la Raspberry Pi, el \gls{Arduino}, el driver de motores, el regulador de tensión, la batería y el sensor RPLidar A1.

\utemFigurePropia{2.Texto/Imag/fig_prototipo_lateral_izquierda.jpg}{0.55\linewidth}
{Vista lateral del prototipo construido.}
{fig:prototipo}

\subsection{Parámetros físicos del modelo}
\label{sec:parametros_modelo}

La Tabla~\ref{tab:parametros_modelo} reúne los parámetros físicos y temporales configurados para el modelo cinemático, la odometría, el gemelo digital y el firmware. Existen tres parámetros nominales de 30~Hz que coinciden numéricamente, pero pertenecen a mecanismos independientes: \texttt{update\_rate} del \texttt{controller\_manager}, \texttt{loop\_rate} de \texttt{DiffDriveArduino} y \texttt{PID\_RATE} del firmware. En el Arduino, la división entera \texttt{1000/30} programa incrementos de 33~ms, según la Ec.~\eqref{eq:pid_intervalo_real}.

\begin{table}[H]
\centering
\caption{Parámetros físicos y temporales empleados en el modelo, la interfaz de hardware y el firmware.}
\label{tab:parametros_modelo}
\small
\begin{tabular}{c p{4.2cm} p{3.0cm} p{4.2cm}}
\toprule
\textbf{Símbolo} & \textbf{Parámetro} & \textbf{Valor} & \textbf{Uso principal} \\
\midrule
$r$ & Radio de rueda & 0.0335 m & Modelo cinemático y odometría \\
$b$ & Separación entre ruedas & 0.188 m & Modelo cinemático diferencial \\
$N_g$ & Relación de reducción & 46 (nominal $46{:}1$) & Cálculo teórico del encoder \\
$N_e^{\mathrm{teo}}$ & Resolución teórica del encoder & 2024 ticks/vuelta & Cálculo teórico \\
$N_e$ & Resolución efectiva del encoder & 1980 ticks/vuelta & Firmware y odometría \\
$c_q$ & Factor de cuadratura & 4 & Decodificación A/B \\
$f_{\mathrm{ctrl}}$ & Frecuencia nominal del ciclo de \texttt{ros2\_control} & 30 Hz & \texttt{controller\_manager.update\_rate} \\
$f_{\mathrm{loop}}$ & Frecuencia nominal usada en la conversión de consignas & 30 Hz & Parámetro \texttt{loop\_rate} de \texttt{DiffDriveArduino} \\
$f_{\mathrm{PID,nom}}$ & Frecuencia nominal declarada del PID embebido & 30 Hz & Constante \texttt{PID\_RATE} del firmware \\
$T_{\mathrm{PID,sched}}$ & Incremento entero del planificador PID & 0.033 s (33 ms) & Actualización de \texttt{nextPID} en el firmware \\
\bottomrule
\end{tabular}
\end{table}

%=====================================================================
\section{Gemelo digital del robot}
\label{ch:robot_digital}

Sobre la base del propósito de validación expuesto en la Sec.~\ref{sec:gemelo_digital}, el gemelo digital del robot se implementa en \gls{ROS2} Foxy mediante URDF/\texttt{xacro} y Gazebo. El modelo principal se encuentra en \path{articubot_one/description/robot.urdf.xacro} e integra los módulos de chasis, ruedas, sensor láser y control. El radio de rueda y la separación entre ruedas de la Tabla~\ref{tab:parametros_modelo} se emplean como referencia geométrica común en el robot físico, el URDF/\texttt{xacro} y \texttt{diff\_drive\_controller}. En cambio, la resolución de encoders y las frecuencias \texttt{loop\_rate}/\texttt{PID\_RATE} pertenecen a la ruta de hardware real: el backend \texttt{GazeboSystem} entrega directamente los estados de las articulaciones y no utiliza la conversión por ticks ni el planificador PID del firmware.

\subsection{Arquitectura del modelo digital}
\label{sec:atribucion}

El modelo digital se organiza mediante archivos \texttt{xacro} modulares dentro de \path{articubot_one/description/}. Desde \texttt{robot.urdf.xacro} se incluye la estructura mecánica definida en \texttt{robot\_core.xacro} y se selecciona el sensor según el entorno: \texttt{lidar.xacro} para hardware real y \texttt{lidar\_sim.xacro} para simulación. La función y localización de cada archivo se consolidan en el \ref{anx:rutas_software}.

La contribución principal en esta etapa corresponde al diseño del modelo URDF/\texttt{xacro}, la parametrización geométrica del robot, la definición de marcos de referencia y la integración con los sistemas de control, SLAM y navegación. La atribución completa de paquetes propios, adaptados y externos se documenta en el \ref{anx:gemelo_digital}.

\subsection{Estructura URDF/xacro y marcos de referencia}
\label{sec:tf-convenciones}

El modelo se organiza como un árbol de eslabones y articulaciones cuyo marco principal es \texttt{base\_link}. En simulación, \texttt{slam\_toolbox} y AMCL utilizan \texttt{base\_footprint}, mientras que en hardware real se utiliza \texttt{base\_link}. La relación exacta entre ambos marcos y la verificación estructural del árbol se documentan en el \ref{anx:gemelo_digital}.

Se adoptan las convenciones REP-103 y REP-105: eje \texttt{x} hacia adelante, eje \texttt{y} hacia la izquierda y eje \texttt{z} hacia arriba \citep{rep103,rep105}. La cadena de transformaciones principal es:

\begin{center}
\texttt{map} $\rightarrow$ \texttt{odom} $\rightarrow$ \texttt{base\_link} $\rightarrow$ \texttt{laser\_frame},
\end{center}

\noindent con la rama fija y coincidente \texttt{base\_link}$\rightarrow$\texttt{base\_footprint}.

\begin{table}[H]
\centering
\caption{Marcos y tópicos principales del modelo digital.}
\label{tab:topics_frames}
\small
\begin{tabular}{p{3.5cm} p{4.0cm} p{5.0cm}}
\toprule
\textbf{Componente} & \textbf{Marco} & \textbf{Tópico asociado} \\
\midrule
Base diferencial & \texttt{odom} $\rightarrow$ \texttt{base\_link} & \texttt{/odom} \\
Sensor láser & \texttt{laser\_frame} & \texttt{/scan} \\
Transformación estática & \texttt{base\_link} $\rightarrow$ \texttt{laser\_frame} & \texttt{/tf\_static} \\
\bottomrule
\end{tabular}
\end{table}

\subsection{Selección del backend de control}
\label{sec:backend}

El modelo permite seleccionar el backend mediante los argumentos \texttt{use\_ros2\_control} y \texttt{sim\_mode}. La ejecución utilizada mantiene \texttt{use\_ros2\_control=true} en ambos entornos: el robot real carga \texttt{diffdrive\_arduino/DiffDriveArduino} y Gazebo carga \texttt{gazebo\_ros2\_control/GazeboSystem}. La trazabilidad del backend alternativo no ejecutado se conserva en el \ref{anx:gemelo_digital}.

\begin{table}[H]
\centering
\caption{Backend de control según entorno de ejecución.}
\label{tab:xacro_args_backend}
\small
\begin{tabular}{p{2.5cm} p{2.5cm} p{8.5cm}}
\toprule
\textbf{Entorno} & \textbf{Backend} & \textbf{Descripción} \\
\midrule
Hardware real & \texttt{ros2\_control} & Plugin \texttt{DiffDriveArduino} y comunicación serial con el \gls{Arduino}. \\
Simulación & \texttt{ros2\_control} & Plugin \texttt{gazebo\_ros2\_control/GazeboSystem} cargado por \texttt{libgazebo\_ros2\_control.so}. \\
\bottomrule
\end{tabular}
\end{table}

\subsection{Sensor láser simulado y criterios de aceptación}

En simulación, el sensor se modela mediante \texttt{lidar\_sim.xacro}, publicando datos en \texttt{/scan} con el marco \texttt{laser\_frame}. El modelo reproduce el comportamiento general del RPLidar A1 físico, aunque con diferencias propias de la simulación, como ausencia de ruido real y un rango mínimo más permisivo; estas diferencias se consideran parte de la brecha simulación--realidad.

Antes de usar el gemelo digital para pruebas de SLAM y navegación, se verifican la continuidad de la cadena TF, la publicación de \texttt{/scan} y \texttt{/odom}, la respuesta cinemática ante comandos de velocidad y la paridad geométrica con el prototipo. La tabla de verificación y la evidencia estructural se trasladan al \ref{anx:gemelo_digital}, Tabla~\ref{tab:criterios_aceptacion_resumen}.

%=====================================================================
\section{Construcción de hardware y comunicación con ROS 2}
\label{chap:robot_real}

Esta sección describe la arquitectura físico-electrónica del robot y su integración con \gls{ROS2} Foxy sobre la Raspberry Pi 4. Los detalles operativos de cableado, reglas \texttt{udev} y configuración extendida se trasladan al \ref{anx:hardware_ros2}.

\subsection{Arquitectura física y de datos}

La plataforma se organiza en dos dominios principales. El dominio de potencia utiliza un bus de 12~V para alimentar los motores DC mediante el driver L298N, y un riel regulado de 5~V para alimentar la Raspberry Pi, el \gls{Arduino}, el \gls{LIDAR} y los periféricos. Esta separación reduce la interferencia eléctrica producida por los motores sobre la electrónica de control.

El dominio de comunicaciones conecta la Raspberry Pi con el \gls{Arduino} mediante enlace serial USB (puerto fijado por regla \texttt{udev} en \texttt{/dev/motor\_controller}, a 57600~baudios), utilizado para enviar referencias de velocidad y recibir datos de encoders. El sensor RPLidar A1 entrega barridos láser publicados en \texttt{/scan}, empleados por \texttt{slam\_toolbox} y por los mapas de costo de Nav2. La Figura~\ref{fig:esquema_conexiones} resume las conexiones eléctricas y de datos del prototipo.

\utemFigurePropia{2.Texto/Imag/circuit_diagram_bb.png}{0.90\textwidth}
{Esquema de conexiones eléctricas y de datos del robot diferencial.}
{fig:esquema_conexiones}

\subsection{Integración con ROS 2}

En hardware real, el lanzamiento \path{articubot_one/launch/launch_robot.launch.py} inicia los nodos necesarios para publicar la descripción del robot, activar \texttt{ros2\_control} y cargar el controlador diferencial. La arquitectura incluye:

\begin{itemize}
\item \texttt{robot\_state\_publisher}, encargado de publicar las transformaciones del robot.
\item \texttt{controller\_manager}, que gestiona la interfaz de hardware y los controladores.
\item \texttt{diff\_drive\_controller}, que convierte comandos de velocidad en referencias para las ruedas.
\item \texttt{joint\_state\_broadcaster}, que publica el estado de las articulaciones.
\end{itemize}

La cadena de transformaciones mantiene la misma convención usada en simulación (Sec.~\ref{sec:tf-convenciones}): \texttt{slam\_toolbox} publica \texttt{map}$\rightarrow$\texttt{odom}, mientras que \texttt{diff\_drive\_controller} publica \texttt{odom}$\rightarrow$\texttt{base\_link} a partir de la odometría de encoders.

\subsection{Teleoperación y multiplexación de comandos}

El nodo \texttt{twist\_mux} arbitra las fuentes de velocidad resumidas en la Tabla~\ref{tab:twist_mux_resumen}. La teleoperación se realiza mediante un mando inalámbrico (nodos \texttt{joy\_node} y \texttt{teleop\_twist\_joy}) y funciona como mecanismo de seguridad mediante un \textit{dead-man switch} asociado al botón de habilitación. La salida del multiplexor se remapea hacia \texttt{/diff\_cont/cmd\_vel\_unstamped}, tópico de entrada del controlador diferencial. La configuración asigna prioridad 100 a la teleoperación y prioridad 10 a la navegación autónoma; la tabla completa se conserva en el \ref{anx:twist_mux_config}, Tabla~\ref{tab:twist_mux_resumen}.

\subsection{Control diferencial e interfaz de hardware}
\label{sec:diff_drive_params}

El controlador diferencial recibe comandos de velocidad lineal y angular, y los convierte en velocidades de rueda mediante el modelo cinemático descrito en el Capítulo~\ref{marco_teorico}. Sus parámetros geométricos (\texttt{wheel\_radius} = $r$ y \texttt{wheel\_separation} = $b$) coinciden con el modelo físico y el gemelo digital, según los valores consolidados en la Tabla~\ref{tab:parametros_modelo}.

\begin{table}[H]
\centering
\caption{Parámetros principales del controlador diferencial.}
\label{tab:diff_drive_params}
\small
\begin{tabular}{llp{5.6cm}}
\toprule
\textbf{Parámetro} & \textbf{Valor} & \textbf{Descripción} \\
\midrule
\texttt{linear.x.min/max\_velocity} & $-0.13/0.13$~m/s & Límites de velocidad lineal \\
\texttt{angular.z.min/max\_velocity} & $-0.35/0.35$~rad/s & Límites de velocidad angular \\
\texttt{update\_rate} & 30~Hz & Frecuencia nominal del \texttt{controller\_manager} \\
\texttt{enable\_odom\_tf} & true & Publica \texttt{odom}$\rightarrow$\texttt{base\_link} \\
\bottomrule
\end{tabular}
\end{table}

La comunicación con el firmware se realiza mediante el plugin \texttt{DiffDriveArduino}, una interfaz de sistema de \texttt{ros2\_control} cuyo ciclo \texttt{read()} $\rightarrow$ \texttt{update()} $\rightarrow$ \texttt{write()} se programa nominalmente con $f_{\mathrm{ctrl}}=30$~Hz mediante \texttt{controller\_manager.update\_rate}.

En \texttt{write()}, la velocidad angular de rueda comandada por \texttt{diff\_drive\_controller} se traduce al formato de \textit{ticks/frame} mediante la Ec.~\eqref{eq:mps_to_ticks}. La división utiliza el parámetro independiente \texttt{cfg\_.loop\_rate}, denotado $f_{\mathrm{loop}}=30$~Hz. Aunque $f_{\mathrm{ctrl}}$ y $f_{\mathrm{loop}}$ coinciden numéricamente, no representan el mismo mecanismo ni comparten una medición temporal.

En \texttt{read()}, la transformación desde los conteos acumulados hasta los estados de rueda sigue la cadena:

\begin{equation}
c_j[k]
\xrightarrow{\text{Ec.~\eqref{eq:posicion_rueda_encoder}}}
\phi_j[k]
\xrightarrow[\Delta t_k]{\text{Ec.~\eqref{eq:w_ruedas}}}
\widehat{\omega}_j[k],
\qquad j\in\{R,L\}.
\label{eq:enc_to_state}
\end{equation}

El intervalo $\Delta t_k$ se define en la Ec.~\eqref{eq:intervalo_diffdrive}, y su secuencia de inicialización se formaliza en el Alg.~\ref{alg:ros2control_cycle}. El plugin exporta una única estimación de velocidad angular por rueda y no calcula otra con el período nominal de 30~Hz.

Las ganancias PID no se transmiten en tiempo de ejecución: residen en el firmware con los valores obtenidos por sintonización PSO y refinamiento experimental (Capítulo~\ref{marco_teorico}), de modo que el plugin solo intercambia referencias de velocidad y lecturas de encoder.

\subsection{Flujo de datos del sistema}
\label{sec:flujo_datos}

El flujo de control entre \gls{ROS2}, el plugin de hardware y el firmware se formaliza en el Algoritmo~\ref{alg:ros2control_cycle}: \texttt{read()} obtiene los conteos y exporta los estados de rueda, \texttt{diff\_drive\_controller} actualiza la odometría y las consignas, y \texttt{write()} transmite referencias enteras de \textit{ticks/frame}. El orden exacto, las validaciones seriales y la inicialización temporal se conservan en el Alg.~\ref{alg:ros2control_cycle} del \ref{anx:implementacion_diffdrive}. El firmware incorpora además una auto-detención de seguridad ante pérdida de comandos, detallada en la Sec.~\ref{sec:bucle_principal}.

%=====================================================================
\section{Firmware de bajo nivel en Arduino Nano}
\label{ch:setup_arduino}

El firmware \texttt{ROSArduinoBridge} parte de un proyecto de código abierto homónimo y se adaptó para este trabajo, incorporando ganancias PID independientes por rueda y la resolución efectiva de encoders del prototipo. Cumple cuatro funciones principales: lectura de encoders, control PID de velocidad por rueda, comunicación serial con \gls{ROS2} mediante \texttt{ros2\_control} y auto-detención de seguridad ante pérdida de comandos. Su consistencia con la capa \gls{ROS2} (Sec.~\ref{chap:robot_real}) depende de los parámetros compartidos reunidos en la Tabla~\ref{tab:anx_trazabilidad_parametros}.

\subsection{Arquitectura modular y constantes físicas}

Los módulos del firmware se encuentran en \path{arduino_firmware/ROSArduinoBridge/}. \texttt{ROSArduinoBridge.ino} coordina la comunicación serial y el ciclo principal; \texttt{commands.h} define los comandos; \texttt{encoder\_driver.h}/\texttt{encoder\_driver.ino} realizan la lectura de encoders; \texttt{motor\_driver.h}/\texttt{motor\_driver.ino} accionan el L298N; y \texttt{diff\_controller.h} implementa el controlador PID. El inventario de rutas se conserva en el \ref{anx:rutas_software}.

El firmware utiliza el diámetro de rueda $D=2r=0.067$~m y la resolución efectiva $N_e$ compartidos con la ruta de hardware real (Tabla~\ref{tab:parametros_modelo}), parametrizados como \texttt{WHEEL\_DIAMETER} y \texttt{ENC\_TICKS\_PER\_REV}. La constante \texttt{PID\_RATE}$=30$ pertenece específicamente al planificador del firmware; interviene en las conversiones y en la división entera \texttt{1000/PID\_RATE}, por lo que \texttt{nextPID} se incrementa en 33~ms, no en $1/30$~s exactos. El backend de Gazebo no utiliza \texttt{ENC\_TICKS\_PER\_REV} ni \texttt{PID\_RATE}. A partir de estas constantes se calcula el factor de conversión entre velocidad lineal en m/s y velocidad de rueda en \textit{ticks/frame} según la Ec.~\ref{eq:mps_to_ticks} ($\approx 313.56$), evaluado en el firmware con $D=2r$.

\subsection{Comunicación serial y lectura de encoders}
\label{sec:serial_firmware}

La comunicación entre la Raspberry Pi y el \gls{Arduino} se realiza mediante un protocolo serial textual a 57600~baud. Durante la operación normal, los comandos principales permiten leer los encoders, enviar referencias de velocidad, cargar las ganancias PID y aplicar PWM directo durante las pruebas en lazo abierto. El protocolo completo se presenta en la Tabla~\ref{tab:protocolo_serial}.

\subsubsection{Obtención experimental de ticks por vuelta}
\label{sec:cpr_experimental}

Los contadores de encoder se midieron durante diez vueltas completas por rueda. Las medias obtenidas fueron 1976~ticks/vuelta para la rueda derecha y 1969~ticks/vuelta para la izquierda; a partir de ellas se adoptó el valor unificado $N_e=1980$~ticks/vuelta, con una diferencia respecto de ambas medias inferior al 0.6\%. Las muestras individuales, el tratamiento de la lectura no atómica de los contadores y el cálculo formal de las medias se conservan en el \ref{anx:pid_data}, Ecs.~\eqref{eq:promedio_cuentas_R} y \eqref{eq:promedio_cuentas_L}.

\subsection{Control PID en el firmware}
\label{sec:motorcontrol_pid}

El controlador de velocidad se ejecuta de forma independiente para cada rueda. La referencia y la medición se expresan en \textit{ticks/frame}. El error por rueda se define según la Ec.~\ref{eq:error_pid}, con la medición entera $m_j^{int}[k]$ en ticks/frame.

La ley implementada corresponde a la forma incremental con derivada sobre la medición definida en el Cap.~\ref{marco_teorico}. El firmware incluye anti-windup condicional, zona muerta para reducir ciclos límite por cuantización y saturación de la salida PWM. Las ganancias definitivas son:

\begin{table}[H]
\centering
\caption{Ganancias PID definitivas codificadas en el firmware.}
\label{tab:ganancias_pid}
\small
\begin{tabular}{lcc}
\toprule
\textbf{Parámetro} & \textbf{Rueda izquierda} & \textbf{Rueda derecha} \\
\midrule
$K_p$ & 16.00 & 16.42 \\
$K_d$ & 20.30 & 20.05 \\
$K_i$ & 0.075 & 0.001 \\
$K_o$ & 50 & 50 \\
\bottomrule
\end{tabular}
\end{table}

Las ganancias $K_{p,j},K_{d,j},K_{i,j}$ se expresan en la escala interna del firmware. El código divide por $K_o$ el numerador completo de la Ec.~\ref{eq:pid_incremento} y después asigna el resultado a una variable \texttt{long}. Debido a ese truncamiento y a los estados enteros, la operación ejecutada no debe reemplazarse por tres ganancias predivididas $K_p/K_o$, $K_d/K_o$ y $K_i/K_o$ como si fueran algebraicamente equivalentes. La salida se satura al intervalo PWM $[u_{\min},u_{\max}]=[-255,255]$ (constante \texttt{MAX\_PWM} del firmware). La convención heredada de \texttt{ROSArduinoBridge} ordena los parámetros como $K_p$, $K_d$, $K_i$, $K_o$.

La función \texttt{resetPID()} no se limita a reiniciar estados: también vuelve a cargar las ganancias predeterminadas mostradas en la Tabla~\ref{tab:ganancias_pid}. Se invoca al reiniciar encoders, antes de aplicar PWM directo y cuando ambos comandos de velocidad son exactamente cero. En consecuencia, cualquier ganancia cargada temporalmente mediante \texttt{u} o \texttt{z} se reemplaza por los valores codificados al producirse cualquiera de esos reinicios. La secuencia exacta de lectura, actualización de estados, saturación y aplicación de las salidas se conserva en el Alg.~\ref{alg:pid_firmware} del \ref{anx:firmware_detalle}.

\subsection{Sintonización del controlador PID asistida por PSO}
\label{sec:sintonizacion_pso}

Las ganancias PID se determinaron mediante el procedimiento híbrido descrito en la Sec.~\ref{sec:sintonizacion}: identificación FODT, optimización por PSO y refinamiento experimental, formalizado en el Alg.~\ref{alg:sintonizacion}. La identificación y la sintonización utilizan enjambres de 50 partículas durante 150 generaciones. Las variables, escalas, cotas y funciones de costo se conservan en la Tabla~\ref{tab:pso_config} del \ref{anx:protocolos_sintonizacion}.

\subsubsection{Captura de respuesta en lazo abierto}
\label{sec:captura_lazo_abierto}

La prueba aplica un escalón PWM nominal de 143 y registra los conteos de encoder cada $T_{cap}=0.1$~s durante 10~s. La velocidad se calcula con el período nominal, mientras el tiempo real del ciclo solo determina el remanente de espera. La ecuación y el procedimiento reproducible se trasladan al \ref{anx:protocolos_sintonizacion}, Ec.~\eqref{eq:vel_lazo_abierto} y Alg.~\ref{alg:captura_lazo_abierto}.

\subsubsection{Identificación y optimización}

La identificación FODT ajusta $(K_m,\tau,L_d)$ por PSO sobre la respuesta al escalón en lazo abierto (Ec.~\ref{eq:fodt_discreto}); luego, ese mismo modelo de primer orden actúa como planta en la simulación del controlador (Ec.~\ref{eq:fodt_lazo_cerrado}). Las diferencias entre ambas etapas —paso de integración, tratamiento del retardo, referencia y tipos numéricos— se detallan en la Sec.~\ref{sec:fodt}.

\subsubsection{Refinamiento experimental}
\label{subsec:refinamiento_iterativo}

Las ganancias obtenidas por PSO se usaron como punto de partida y luego se ajustaron sobre el robot real (Etapa~3 del Alg.~\ref{alg:sintonizacion}). El código final solo conserva las ganancias resultantes por rueda; no registra de forma automática qué modificación experimental se asocia a cada fenómeno mecánico.

La salida de la ejecución utilizada como punto de partida, incluidos los parámetros FODT identificados y las ganancias sugeridas por PSO, se conserva en la Tabla~\ref{tab:pso_resultados} del \ref{anx:resultados_pso_ejecucion}. Al no fijarse una semilla pseudoaleatoria, los valores pueden variar entre ejecuciones. Además, con $T_{cap}=0.1$~s, los retardos obtenidos en esa ejecución son inferiores a un período de captura y no producen un desplazamiento temporal efectivo en el modelo discreto.

\subsection{Bucle principal y auto-detención}
\label{sec:bucle_principal}

El bucle principal atiende el parser serial, programa el PID mediante incrementos enteros de 33~ms y evalúa la auto-detención. Si transcurren más de 2000~ms sin recibir comandos de movimiento, el firmware detiene los motores y desactiva la bandera de movimiento. La rama de auto-detención no llama directamente a \texttt{resetPID()}; en la siguiente actualización programada, esta función solo se ejecuta si alguno de los valores \texttt{PrevInput} es distinto de cero. Esta condición evita que el robot continúe desplazándose ante pérdida de conexión o falla de algún nodo de navegación.

%=====================================================================
\section{Arquitectura de control y navegación: SLAM, Nav2 y exploración}
\label{chap:app_control_nav}

Las capas de hardware, firmware, gemelo digital y control de bajo nivel fueron desarrolladas en las secciones anteriores; esta sección enfatiza la orquestación de los subsistemas mediante \texttt{ros2\_control}, \texttt{slam\_toolbox}, Nav2 y \texttt{explore\_lite}. Las configuraciones ejecutables se mantienen en \path{articubot_one/config/}; la correspondencia entre archivos, entorno y función se reúne en el \ref{anx:rutas_software}, y los parámetros completos de configuración se trasladan al \ref{anx:nav2_slam_config}.

\subsection{Mapeo con \texttt{slam\_toolbox}}
\label{sec:slam_ch9}

El robot construye el mapa del entorno mediante \texttt{slam\_toolbox} (fundamentos en Sec.~\ref{sec:slam}), en modo \texttt{online\_async}, que procesa los barridos de forma asíncrona. En la arquitectura distribuida ejecutada, este nodo se despliega en la estación de trabajo; la Raspberry Pi 4 publica \texttt{/scan}, \texttt{/odom} y las transformaciones asociadas al robot. La optimización del grafo permite corregir la deriva acumulada durante el recorrido.

La configuración real utiliza Ceres, los marcos \texttt{map}/\texttt{odom}/\texttt{base\_link}, el tópico \texttt{/scan} y cierre de lazo. Los valores completos se trasladan a la Tabla~\ref{tab:slam_params_ch9} del \ref{anx:nav2_slam_config}.

\subsection{Navegación autónoma con Nav2}
\label{sec:nav2_ch9}

Nav2 recibe una meta de navegación, calcula una ruta global y genera comandos de velocidad para que el robot la ejecute evitando obstáculos. La cadena funcional implementada es:

\begin{center}
\small
\begin{minipage}{\linewidth}
\centering
Meta $\rightarrow$ \texttt{bt\_navigator} $\rightarrow$ \texttt{planner\_server} $\rightarrow$ \\[3pt]
\texttt{controller\_server} $\rightarrow$ \texttt{twist\_mux} $\rightarrow$ \texttt{diff\_drive\_controller}
\end{minipage}
\end{center}

El planificador global utilizado es NavFn con A$^*$ (Sec.~\ref{sec:nav_marco}), activado mediante \texttt{use\_astar} sobre el \texttt{global\_costmap}. Para el seguimiento local se emplea DWB, que evalúa trayectorias candidatas y selecciona comandos de velocidad compatibles con las restricciones cinemáticas del robot.

\begin{table}[H]
\centering
\caption{Componentes principales de Nav2 utilizados.}
\label{tab:nav2_componentes_ch9}
\small
\begin{tabular}{llp{6.5cm}}
\toprule
\textbf{Componente} & \textbf{Implementación} & \textbf{Función} \\
\midrule
Planificador global & NavFn con A$^*$ & Generar ruta global \\
Controlador local & DWB & Seguir ruta y evitar obstáculos \\
Costmaps & Global y local & Representar obstáculos y zonas transitables \\
Árbol de comportamiento & \texttt{bt\_navigator} & Orquestar navegación y recuperación \\
Multiplexor & \texttt{twist\_mux} & Priorizar navegación o teleoperación \\
\bottomrule
\end{tabular}
\end{table}

\subsection{Mapas de costos}
\label{sec:costmaps_ch9}

Nav2 utiliza dos mapas de costos. El \texttt{global\_costmap} entrega la representación usada por el planificador global, mientras que el \texttt{local\_costmap} se actualiza alrededor del robot y permite al controlador DWB reaccionar ante obstáculos cercanos. Ambos mapas incorporan una capa de inflación, necesaria para mantener una distancia de seguridad respecto de los obstáculos.

En hardware real, las capas de inflación de \texttt{global\_costmap} y \texttt{local\_costmap} utilizan \texttt{inflation\_radius}=0.15~m y \texttt{cost\_scaling\_factor}=5.0, con resolución de 0.05~m/celda. En simulación, ambos costmaps utilizan 0.55~m y 3.0, respectivamente. El detalle completo se presenta en el \ref{anx:nav2_slam_config}.

\subsection{Árbol de comportamiento y recuperación}
\label{sec:bt_ch9}

La navegación de alto nivel se organiza mediante un árbol de comportamiento. La rama principal calcula la ruta global y ejecuta el seguimiento local; si ocurre un fallo, se activan acciones de recuperación, lo que permite reintentos ante bloqueos sin modificar la lógica completa del sistema.

Las acciones configuradas (Tabla~\ref{tab:recovery_ch9}) siguen un principio de escalada: primero acciones de bajo costo computacional (limpieza de costmaps) y luego maniobras físicas (retroceso, giro y espera), con replanificación entre reintentos. El detalle de cada acción se conserva en el \ref{anx:recovery_nav2}, Tabla~\ref{tab:recovery_ch9}.

\subsection{Exploración autónoma de fronteras}
\label{sec:explore_ch9}

Para mapear entornos desconocidos sin intervención manual continua, se integra el módulo de exploración por fronteras seleccionado en la Sec.~\ref{sec:exploracion} (\texttt{explore\_lite}); su nodo \texttt{explore\_server} identifica, sobre el costmap global, las fronteras entre espacio libre conocido y zonas no observadas, y envía metas de navegación a Nav2. De esta forma, se cierra el ciclo percepción--mapeo--navegación--exploración sobre la misma arquitectura.

El explorador utiliza el costmap global, reevalúa las fronteras con una frecuencia configurada de 0.3~Hz y aplica un tiempo máximo sin progreso de 10~s. Los parámetros completos se conservan en la Tabla~\ref{tab:explore_params_ch9} del \ref{anx:nav2_slam_config}, y el flujo reproducible de detección, selección, envío de metas y gestión de la lista negra se traslada al Alg.~\ref{alg:exploracion_fronteras} del mismo anexo.

\subsection{Integración del ciclo autónomo y transferencia simulación--hardware}
\label{sec:integracion_ciclo}

Los subsistemas operan de forma encadenada: \texttt{slam\_toolbox} construye el mapa y publica la transformación \texttt{map}$\rightarrow$\texttt{odom}; Nav2 utiliza ese mapa para planificar y controlar el movimiento; \texttt{twist\_mux} envía los comandos al controlador diferencial; y \texttt{explore\_server} genera nuevas metas hacia fronteras no exploradas. Esta integración permite que el robot pase desde teleoperación y mapeo inicial hacia navegación autónoma y exploración de zonas desconocidas.

\label{sec:verificacion_ch9}
La transferencia al hardware se realiza manteniendo la estructura de tópicos y marcos, pero seleccionando archivos de configuración específicos para el entorno real; la asignación de \texttt{use\_sim\_time} y los archivos YAML equivalentes de simulación y hardware se detallan en el \ref{anx:nav2_slam_config}. Además, se ajustan límites cinemáticos, frecuencias, tolerancias e inflación de costmaps según las capacidades del prototipo físico.

%=====================================================================
\section*{Conclusión del capítulo}
\label{sec:conclusion_ch9}

La implementación integra componentes accesibles —plataforma diferencial, encoders, \gls{Arduino} Nano, Raspberry Pi 4 y RPLidar A1— con una arquitectura de software basada en \gls{ROS2}. El gemelo digital, construido sobre los mismos parámetros geométricos del prototipo, permitió validar de forma previa la cinemática, los marcos de referencia y la configuración de control, facilitando la transferencia hacia el hardware real. La comunicación mediante \texttt{ros2\_control}, \texttt{diff\_drive\_controller} y el plugin \texttt{DiffDriveArduino} cierra el ciclo entre comandos de navegación, control PID de ruedas —sintonizado mediante PSO y refinamiento experimental— y publicación de odometría, con auto-detención de seguridad ante pérdida de comandos. Sobre esta base, la orquestación de \texttt{slam\_toolbox}, Nav2 y \texttt{explore\_lite} articula percepción, mapeo, navegación y exploración autónoma sobre una misma estructura de tópicos y marcos, tanto en simulación como en hardware real. Su desempeño se evalúa en el Capítulo~\ref{ch:resultados}.
